\Chapter{72}

\Verse{1}
{
    \zA{且说鸳鸯出了角门，脸上犹热，心内突突的乱跳，真是意外之事。}
}
{
    \eA{[English source not present in the checked-in Bencraft/Joly files.]}
}
{
}

\Verse{2}
{
    \zA{因想这事非 常，若说出来奸盗相连，关系人命，还保不住带累旁人。}
}
{
    \eA{[English source not present in the checked-in Bencraft/Joly files.]}
}
{
}

\Verse{3}
{
    \zA{横竖与自己无干，且藏在 心内，不说给人知道。}
}
{
    \eA{[English source not present in the checked-in Bencraft/Joly files.]}
}
{
}

\Verse{4}
{
    \zA{回房复了贾母的命，大家安息不提。}
}
{
    \eA{[English source not present in the checked-in Bencraft/Joly files.]}
}
{
}

\Verse{5}
{
    \zA{却说司棋因从小儿和他姑表兄弟一处玩笑，起初时小儿戏言，便都订下将来不 娶不嫁；近年大了，彼此又出落得品貌风流。}
}
{
    \eA{[English source not present in the checked-in Bencraft/Joly files.]}
}
{
}

\Verse{6}
{
    \zA{常时司棋回家时，二人眉来眼去，旧 情不断，只不能入手。}
}
{
    \eA{[English source not present in the checked-in Bencraft/Joly files.]}
}
{
}

\Verse{7}
{
    \zA{又彼此生怕父母不从，二人便设法，彼此里外买嘱园内老婆 子们，留门看道。}
}
{
    \eA{[English source not present in the checked-in Bencraft/Joly files.]}
}
{
}

\Verse{8}
{
    \zA{今日赶乱，方从外进来，初次入港。}
}
{
    \eA{[English source not present in the checked-in Bencraft/Joly files.]}
}
{
}

\Verse{9}
{
    \zA{虽未成双，却也海誓山盟， 私传表记，已有无限风情。}
}
{
    \eA{[English source not present in the checked-in Bencraft/Joly files.]}
}
{
}

\Verse{10}
{
    \zA{忽被鸳鸯惊散，那小厮早穿花度柳，从角门出去了。}
}
{
    \eA{[English source not present in the checked-in Bencraft/Joly files.]}
}
{
}

\Verse{11}
{
    \zA{司 棋一夜不曾睡着，又后悔不来。}
}
{
    \eA{[English source not present in the checked-in Bencraft/Joly files.]}
}
{
}

\Verse{12}
{
    \zA{至次日见了鸳鸯，自是脸上一红一白，百般过不去， 心内怀着鬼胎，茶饭无心，起坐恍惚。}
}
{
    \eA{[English source not present in the checked-in Bencraft/Joly files.]}
}
{
}

\Verse{13}
{
    \zA{挨了两日，竟不听见有动静，方略放下了心。}
}
{
    \eA{[English source not present in the checked-in Bencraft/Joly files.]}
}
{
}

\Verse{14}
{
    \zA{这日晚间，忽有个婆子来悄悄告诉道：“你表兄竟逃走了，三四天没上家。}
}
{
    \eA{[English source not present in the checked-in Bencraft/Joly files.]}
}
{
}

\Verse{15}
{
    \zA{如今打 发人四处找他呢。”}
}
{
    \eA{[English source not present in the checked-in Bencraft/Joly files.]}
}
{
}

\Verse{16}
{
    \zA{司棋听了，又急又气又伤心，因想道：“纵然闹出来，也该死在 一处。}
}
{
    \eA{[English source not present in the checked-in Bencraft/Joly files.]}
}
{
}

\Verse{17}
{
    \zA{真真男人没情意，先就走了。”}
}
{
    \eA{[English source not present in the checked-in Bencraft/Joly files.]}
}
{
}

\Verse{18}
{
    \zA{因此，又添了一层气，次日便觉心内不快， 支持不住，一头躺倒，恹恹的成了病了。}
}
{
    \eA{[English source not present in the checked-in Bencraft/Joly files.]}
}
{
}

\Verse{19}
{
    \zA{鸳鸯闻知那边无故走了一个小厮，园内司棋病重，要往外挪，心下料定是二人 惧罪之故，“生怕我说出来。”}
}
{
    \eA{[English source not present in the checked-in Bencraft/Joly files.]}
}
{
}

\Verse{20}
{
    \zA{因此，自己反过意不去，指着来望候司棋，支出人去， 反自己赌咒发誓，与司棋说：“我若告诉一个人，立刻现死现报!你只管放心养病， 别白遭塌了小命儿。”}
}
{
    \eA{[English source not present in the checked-in Bencraft/Joly files.]}
}
{
}

\Verse{21}
{
    \zA{司棋一把拉住，哭道：“我的姐姐!咱们从小儿耳鬓厮磨，你 不曾拿我当外人待，我也不敢怠慢了你，如今我虽一着走错了，你若果然不告诉一
        个人，你就是我的亲娘一样。}
}
{
    \eA{[English source not present in the checked-in Bencraft/Joly files.]}
}
{
}

\Verse{22}
{
    \zA{从此后，我活一日，是你给我一日。}
}
{
    \eA{[English source not present in the checked-in Bencraft/Joly files.]}
}
{
}

\Verse{23}
{
    \zA{我的病要好了， 把你立个长生牌位，我天天烧香磕头，保佑你一辈子福寿双全的，我若死了时，变 驴变狗报答你。}
}
{
    \eA{[English source not present in the checked-in Bencraft/Joly files.]}
}
{
}

\Verse{24}
{
    \zA{倘或咱们散了，以后遇见，我自有报答的去处。”}
}
{
    \eA{[English source not present in the checked-in Bencraft/Joly files.]}
}
{
}

\Verse{25}
{
    \zA{一面说，一面哭。}
}
{
    \eA{[English source not present in the checked-in Bencraft/Joly files.]}
}
{
}

\Verse{26}
{
    \zA{这一席话，反把鸳鸯说的酸心，也哭起来了。}
}
{
    \eA{[English source not present in the checked-in Bencraft/Joly files.]}
}
{
}

\Verse{27}
{
    \zA{因点头道：“你也是自家要作死哟!我 作什么管你这些事坏你的名儿，我白去献勤儿？}
}
{
    \eA{[English source not present in the checked-in Bencraft/Joly files.]}
}
{
}

\Verse{28}
{
    \zA{况且这事我也不便开口和人说。}
}
{
    \eA{[English source not present in the checked-in Bencraft/Joly files.]}
}
{
}

\Verse{29}
{
    \zA{你 只放心。}
}
{
    \eA{[English source not present in the checked-in Bencraft/Joly files.]}
}
{
}

\Verse{30}
{
    \zA{从此养好了，可要安分守己的，再别胡行乱闹了。”}
}
{
    \eA{[English source not present in the checked-in Bencraft/Joly files.]}
}
{
}

\Verse{31}
{
    \zA{司棋在枕上点首不绝。}
}
{
    \eA{[English source not present in the checked-in Bencraft/Joly files.]}
}
{
}

\Verse{32}
{
    \zA{鸳鸯又安慰了他一番，方出来。}
}
{
    \eA{[English source not present in the checked-in Bencraft/Joly files.]}
}
{
}

\Verse{33}
{
    \zA{因知贾琏不在家中，又因这两日凤姐儿声色怠 惰了些，不似往日一样，便顺路来问候。}
}
{
    \eA{[English source not present in the checked-in Bencraft/Joly files.]}
}
{
}

\Verse{34}
{
    \zA{刚进入凤姐院中，二门上的人见是他来， 便站立待他进去。}
}
{
    \eA{[English source not present in the checked-in Bencraft/Joly files.]}
}
{
}

\Verse{35}
{
    \zA{鸳鸯来至堂屋，只见平儿从里头出来，见了他来，便忙上来悄声 笑道：“才吃了一口饭，歇了中觉了。}
}
{
    \eA{[English source not present in the checked-in Bencraft/Joly files.]}
}
{
}

\Verse{36}
{
    \zA{你且这屋里略坐坐。”}
}
{
    \eA{[English source not present in the checked-in Bencraft/Joly files.]}
}
{
}

\Verse{37}
{
    \zA{鸳鸯听了，只得同平儿 到东边房里来。}
}
{
    \eA{[English source not present in the checked-in Bencraft/Joly files.]}
}
{
}

\Verse{38}
{
    \zA{小丫头倒了茶来。}
}
{
    \eA{[English source not present in the checked-in Bencraft/Joly files.]}
}
{
}

\Verse{39}
{
    \zA{鸳鸯悄问道：“你奶奶这两日是怎么了？}
}
{
    \eA{[English source not present in the checked-in Bencraft/Joly files.]}
}
{
}

\Verse{40}
{
    \zA{我近来看 着他懒懒的。”}
}
{
    \eA{[English source not present in the checked-in Bencraft/Joly files.]}
}
{
}

\Verse{41}
{
    \zA{平儿见问，因房内无人，便叹道：“他这懒懒的，也不止今日了。}
}
{
    \eA{[English source not present in the checked-in Bencraft/Joly files.]}
}
{
}

\Verse{42}
{
    \zA{这 有一月前头，就是这么着。}
}
{
    \eA{[English source not present in the checked-in Bencraft/Joly files.]}
}
{
}

\Verse{43}
{
    \zA{这几日忙乱了几天，又受了些闲气，从新又勾起来。}
}
{
    \eA{[English source not present in the checked-in Bencraft/Joly files.]}
}
{
}

\Verse{44}
{
    \zA{这 两日比先又添了些病，所以支不住，就露出马脚来了。”}
}
{
    \eA{[English source not present in the checked-in Bencraft/Joly files.]}
}
{
}

\Verse{45}
{
    \zA{鸳鸯道：“既这样，怎么不 早请大夫治？”}
}
{
    \eA{[English source not present in the checked-in Bencraft/Joly files.]}
}
{
}

\Verse{46}
{
    \zA{平儿叹道：“我的姐姐，你还不知道他那脾气的？}
}
{
    \eA{[English source not present in the checked-in Bencraft/Joly files.]}
}
{
}

\Verse{47}
{
    \zA{别说请大夫来吃药， 我看不过，白问一声‘身上觉怎么样’，他就动了气，反说我咒他病了。}
}
{
    \eA{[English source not present in the checked-in Bencraft/Joly files.]}
}
{
}

\Verse{48}
{
    \zA{饶这样， 天天还是察三访四。}
}
{
    \eA{[English source not present in the checked-in Bencraft/Joly files.]}
}
{
}

\Verse{49}
{
    \zA{自己再不看破些，且养身子！”}
}
{
    \eA{[English source not present in the checked-in Bencraft/Joly files.]}
}
{
}

\Verse{50}
{
    \zA{鸳鸯道：“虽然如此，到底该请 大夫来瞧瞧是什么病，也都好放心。”}
}
{
    \eA{[English source not present in the checked-in Bencraft/Joly files.]}
}
{
}

\Verse{51}
{
    \zA{平儿叹道：“说起病来，据我看也不是什么小 症候。”}
}
{
    \eA{[English source not present in the checked-in Bencraft/Joly files.]}
}
{
}

\Verse{52}
{
    \zA{鸳鸯忙道：“是什么病呢？”}
}
{
    \eA{[English source not present in the checked-in Bencraft/Joly files.]}
}
{
}

\Verse{53}
{
    \zA{平儿见问，又往前凑了一凑，向耳边说道：“只 从上月行了经之后，这一个月，竟沥沥淅淅的没有止住。}
}
{
    \eA{[English source not present in the checked-in Bencraft/Joly files.]}
}
{
}

\Verse{54}
{
    \zA{这可是大病不是？”}
}
{
    \eA{[English source not present in the checked-in Bencraft/Joly files.]}
}
{
}

\Verse{55}
{
    \zA{鸳鸯 听了忙答应道：“嗳哟，依这么说，可不成了‘血山崩’了吗？”}
}
{
    \eA{[English source not present in the checked-in Bencraft/Joly files.]}
}
{
}

\Verse{56}
{
    \zA{平儿忙啐了一口， 又悄笑道：“你个女孩儿家，这是怎么说？}
}
{
    \eA{[English source not present in the checked-in Bencraft/Joly files.]}
}
{
}

\Verse{57}
{
    \zA{你倒会咒人。”}
}
{
    \eA{[English source not present in the checked-in Bencraft/Joly files.]}
}
{
}

\Verse{58}
{
    \zA{鸳鸯见说，不禁红了脸， 又悄笑道：“究竟我也不懂什么是崩不崩的。}
}
{
    \eA{[English source not present in the checked-in Bencraft/Joly files.]}
}
{
}

\Verse{59}
{
    \zA{你倒忘了不成：先我姐姐不是害这病 死了？}
}
{
    \eA{[English source not present in the checked-in Bencraft/Joly files.]}
}
{
}

\Verse{60}
{
    \zA{我也不知是什么病，因无心中听见妈和亲家妈说，我还纳闷，后来听见原故， 才明白了一二分。”}
}
{
    \eA{[English source not present in the checked-in Bencraft/Joly files.]}
}
{
}

\Verse{61}
{
    \zA{二人正说着，只见小丫头向平儿道：“方才朱大娘又来了。}
}
{
    \eA{[English source not present in the checked-in Bencraft/Joly files.]}
}
{
}

\Verse{62}
{
    \zA{我们 回了他：‘奶奶才歇中觉。’}
}
{
    \eA{[English source not present in the checked-in Bencraft/Joly files.]}
}
{
}

\Verse{63}
{
    \zA{他往太太上头去了。”}
}
{
    \eA{[English source not present in the checked-in Bencraft/Joly files.]}
}
{
}

\Verse{64}
{
    \zA{平儿听了点头。}
}
{
    \eA{[English source not present in the checked-in Bencraft/Joly files.]}
}
{
}

\Verse{65}
{
    \zA{鸳鸯问：“那一个 朱大娘？”}
}
{
    \eA{[English source not present in the checked-in Bencraft/Joly files.]}
}
{
}

\Verse{66}
{
    \zA{平儿道：“就是官媒婆朱嫂子。}
}
{
    \eA{[English source not present in the checked-in Bencraft/Joly files.]}
}
{
}

\Verse{67}
{
    \zA{因有个什么孙大人来和咱们求亲，所以 他这两日天天弄个帖子来，闹得人怪烦的。”}
}
{
    \eA{[English source not present in the checked-in Bencraft/Joly files.]}
}
{
}

\Verse{68}
{
    \zA{一语未了，小丫头跑来说：“二爷进来了。”}
}
{
    \eA{[English source not present in the checked-in Bencraft/Joly files.]}
}
{
}

\Verse{69}
{
    \zA{说话之间，贾琏已走至堂屋门口， 平儿忙迎出来。}
}
{
    \eA{[English source not present in the checked-in Bencraft/Joly files.]}
}
{
}

\Verse{70}
{
    \zA{贾琏见平儿在东屋里，便也过这间房内来，走至门前，忽见鸳鸯坐 在炕上，便煞住脚，笑道：“鸳鸯姐姐，今儿贵步幸临贱地！”}
}
{
    \eA{[English source not present in the checked-in Bencraft/Joly files.]}
}
{
}

\Verse{71}
{
    \zA{鸳鸯只坐着，笑道： “来请爷奶奶的安，偏又不在家的不在家，睡觉的睡觉。”}
}
{
    \eA{[English source not present in the checked-in Bencraft/Joly files.]}
}
{
}

\Verse{72}
{
    \zA{贾琏笑道：“姐姐一年到 头辛苦，伏侍老太太，我还没看你去，那里还敢劳动来看我们。”}
}
{
    \eA{[English source not present in the checked-in Bencraft/Joly files.]}
}
{
}

\Verse{73}
{
    \zA{又说：“巧的很。}
}
{
    \eA{[English source not present in the checked-in Bencraft/Joly files.]}
}
{
}

\Verse{74}
{
    \zA{我才要找姐姐去，因为穿着这袍子热，先来换了夹袍子，再过去找姐姐去，不想老 天爷可怜，省我走这一趟。”}
}
{
    \eA{[English source not present in the checked-in Bencraft/Joly files.]}
}
{
}

\Verse{75}
{
    \zA{一面说，一面在椅子上坐下。}
}
{
    \eA{[English source not present in the checked-in Bencraft/Joly files.]}
}
{
}

\Verse{76}
{
    \zA{鸳鸯因问：“又有什么说 的？”}
}
{
    \eA{[English source not present in the checked-in Bencraft/Joly files.]}
}
{
}

\Verse{77}
{
    \zA{贾琏未语先笑，道：“因有一件事竟忘了，只怕姐姐还记得：上年老太太生 日，曾有一个外路和尚来孝敬一个腊油冻的佛手，因老太太爱，就即刻拿过来摆着。}
}
{
    \eA{[English source not present in the checked-in Bencraft/Joly files.]}
}
{
}

\Verse{78}
{
    \zA{因前日老太太的生日，我看古董帐，还有一笔在这帐上，却不知此时这件着落在何 处。}
}
{
    \eA{[English source not present in the checked-in Bencraft/Joly files.]}
}
{
}

\Verse{79}
{
    \zA{古董房里的人也回过了我两次，等我问准了，好注了一笔。}
}
{
    \eA{[English source not present in the checked-in Bencraft/Joly files.]}
}
{
}

\Verse{80}
{
    \zA{所以我问姐姐：如 今还是老太太摆着呢，还是交到谁手里去了呢？”}
}
{
    \eA{[English source not present in the checked-in Bencraft/Joly files.]}
}
{
}

\Verse{81}
{
    \zA{鸳鸯听说，便说道：“老太太摆 了几日，厌烦了，就给你们奶奶了，你这会子又问我来了。}
}
{
    \eA{[English source not present in the checked-in Bencraft/Joly files.]}
}
{
}

\Verse{82}
{
    \zA{我连日子还记得，还是 我打发了老王家的送来。}
}
{
    \eA{[English source not present in the checked-in Bencraft/Joly files.]}
}
{
}

\Verse{83}
{
    \zA{你忘了，或是问你们奶奶和平儿。”}
}
{
    \eA{[English source not present in the checked-in Bencraft/Joly files.]}
}
{
}

\Verse{84}
{
    \zA{平儿正拿衣裳，听见 如此说，忙出来回说：“交过来了，现在楼上放着呢。}
}
{
    \eA{[English source not present in the checked-in Bencraft/Joly files.]}
}
{
}

\Verse{85}
{
    \zA{奶奶已经打发人去说过，他 们发昏没记上，又来叨蹬这些没要紧的事。”}
}
{
    \eA{[English source not present in the checked-in Bencraft/Joly files.]}
}
{
}

\Verse{86}
{
    \zA{贾琏听说，笑道：“既然给了你奶奶， 我怎么不知道，你们就昧下了？”}
}
{
    \eA{[English source not present in the checked-in Bencraft/Joly files.]}
}
{
}

\Verse{87}
{
    \zA{平儿道：“奶奶告诉二爷，二爷还要送人，奶奶 不肯，好容易留下的。}
}
{
    \eA{[English source not present in the checked-in Bencraft/Joly files.]}
}
{
}

\Verse{88}
{
    \zA{这会子自己忘了，倒说我们昧下!那是什么好东西？}
}
{
    \eA{[English source not present in the checked-in Bencraft/Joly files.]}
}
{
}

\Verse{89}
{
    \zA{比那强十 倍的也没昧下一遭儿，这会子就爱上那不值钱的咧？”}
}
{
    \eA{[English source not present in the checked-in Bencraft/Joly files.]}
}
{
}

\Verse{90}
{
    \zA{贾琏垂头含笑想了想，拍手 道：“我如今竟糊涂了!丢三忘四，惹人抱怨，竟大不像先了。”}
}
{
    \eA{[English source not present in the checked-in Bencraft/Joly files.]}
}
{
}

\Verse{91}
{
    \zA{鸳鸯笑道：“也怨不 得：事情又多，口舌又杂，你再喝上两钟酒，那里记得许多？”}
}
{
    \eA{[English source not present in the checked-in Bencraft/Joly files.]}
}
{
}

\Verse{92}
{
    \zA{一面说，一面起身 要走。}
}
{
    \eA{[English source not present in the checked-in Bencraft/Joly files.]}
}
{
}

\Verse{93}
{
    \zA{贾琏忙也立起身来，说道：“好姐姐，略坐一坐儿，兄弟还有一事相求。”}
}
{
    \eA{[English source not present in the checked-in Bencraft/Joly files.]}
}
{
}

\Verse{94}
{
    \zA{说着， 便骂小丫头：“怎么不沏好茶来？}
}
{
    \eA{[English source not present in the checked-in Bencraft/Joly files.]}
}
{
}

\Verse{95}
{
    \zA{快拿干净盖碗，把昨日进上的新茶沏一碗来！”}
}
{
    \eA{[English source not present in the checked-in Bencraft/Joly files.]}
}
{
}

\Verse{96}
{
    \zA{说 着，向鸳鸯道：“这两日，因老太太千秋，所有的几千两都使了。}
}
{
    \eA{[English source not present in the checked-in Bencraft/Joly files.]}
}
{
}

\Verse{97}
{
    \zA{几处房租、地租， 统在九月才得，这会子竟接不上。}
}
{
    \eA{[English source not present in the checked-in Bencraft/Joly files.]}
}
{
}

\Verse{98}
{
    \zA{明儿又要送南安府里的礼，又要预备娘娘的重阳 节，还有几家红白大礼，至少还得三二千两银子用，一时难去支借。}
}
{
    \eA{[English source not present in the checked-in Bencraft/Joly files.]}
}
{
}

\Verse{99}
{
    \zA{俗语说的好： ‘求人不如求己。’}
}
{
    \eA{[English source not present in the checked-in Bencraft/Joly files.]}
}
{
}

\Verse{100}
{
    \zA{说不得姐姐担个不是，暂且把老太太查不着的金银家伙，偷着 运出一箱子来，暂押千数两银子，支腾过去。}
}
{
    \eA{[English source not present in the checked-in Bencraft/Joly files.]}
}
{
}

\Verse{101}
{
    \zA{不上半月的光景银子来了，我就赎了 交还，断不能叫姐姐落不是。”}
}
{
    \eA{[English source not present in the checked-in Bencraft/Joly files.]}
}
{
}

\Verse{102}
{
    \zA{鸳鸯听了，笑道：“你倒会变法儿!亏你怎么想了。”}
}
{
    \eA{[English source not present in the checked-in Bencraft/Joly files.]}
}
{
}

\Verse{103}
{
    \zA{贾琏笑道：“不是我撒谎：若论除了姐姐，也还有人手里管得起千数两银子；只是 他们为人都不如你明白有胆量，我和他们一说，反吓住了他们。}
}
{
    \eA{[English source not present in the checked-in Bencraft/Joly files.]}
}
{
}

\Verse{104}
{
    \zA{所以我‘宁撞金钟 一下，不打铙钹三千’。”}
}
{
    \eA{[English source not present in the checked-in Bencraft/Joly files.]}
}
{
}

\Verse{105}
{
    \zA{一语未了，贾母那边小丫头子忙忙走来找鸳鸯，说：“老 太太找姐姐呢。}
}
{
    \eA{[English source not present in the checked-in Bencraft/Joly files.]}
}
{
}

\Verse{106}
{
    \zA{这半日，我那里没找到？}
}
{
    \eA{[English source not present in the checked-in Bencraft/Joly files.]}
}
{
}

\Verse{107}
{
    \zA{却在这里。”}
}
{
    \eA{[English source not present in the checked-in Bencraft/Joly files.]}
}
{
}

\Verse{108}
{
    \zA{鸳鸯听说，忙着去见贾母。}
}
{
    \eA{[English source not present in the checked-in Bencraft/Joly files.]}
}
{
}

\Verse{109}
{
    \zA{贾琏见他去了，只得回来瞧凤姐。}
}
{
    \eA{[English source not present in the checked-in Bencraft/Joly files.]}
}
{
}

\Verse{110}
{
    \zA{谁知凤姐已醒了，听他和鸳鸯借当，自己不 便答话，只躺在榻上。}
}
{
    \eA{[English source not present in the checked-in Bencraft/Joly files.]}
}
{
}

\Verse{111}
{
    \zA{听见鸳鸯去了，贾琏进来，凤姐因问道：“他可应准了？”}
}
{
    \eA{[English source not present in the checked-in Bencraft/Joly files.]}
}
{
}

\Verse{112}
{
    \zA{贾琏笑道：“虽未应准，却有几分成了。}
}
{
    \eA{[English source not present in the checked-in Bencraft/Joly files.]}
}
{
}

\Verse{113}
{
    \zA{须得你再去和他说一说，就十分成了。”}
}
{
    \eA{[English source not present in the checked-in Bencraft/Joly files.]}
}
{
}

\Verse{114}
{
    \zA{凤 姐笑道：“我不管这些事。}
}
{
    \eA{[English source not present in the checked-in Bencraft/Joly files.]}
}
{
}

\Verse{115}
{
    \zA{倘或说准了，这会子说着好听，到了有钱的时节，你就 撂在脖子后头了，谁和你打饥荒去？}
}
{
    \eA{[English source not present in the checked-in Bencraft/Joly files.]}
}
{
}

\Verse{116}
{
    \zA{倘或老太太知道了，倒把我这几年的脸面都丢 了。”}
}
{
    \eA{[English source not present in the checked-in Bencraft/Joly files.]}
}
{
}

\Verse{117}
{
    \zA{贾琏笑道：“好人，你要说定了，我谢你。”}
}
{
    \eA{[English source not present in the checked-in Bencraft/Joly files.]}
}
{
}

\Verse{118}
{
    \zA{凤姐笑道：“你说谢我什么？”}
}
{
    \eA{[English source not present in the checked-in Bencraft/Joly files.]}
}
{
}

\Verse{119}
{
    \zA{贾 琏笑道：“你说要什么就有什么。”}
}
{
    \eA{[English source not present in the checked-in Bencraft/Joly files.]}
}
{
}

\Verse{120}
{
    \zA{平儿一旁笑道：“奶奶不用要别的。}
}
{
    \eA{[English source not present in the checked-in Bencraft/Joly files.]}
}
{
}

\Verse{121}
{
    \zA{刚才正说要 做一件什么事，恰少一二百银子使，不如借了来，奶奶拿这么一二百银子，岂不两 全其美？”}
}
{
    \eA{[English source not present in the checked-in Bencraft/Joly files.]}
}
{
}

\Verse{122}
{
    \zA{凤姐笑道：“幸亏提起我来。}
}
{
    \eA{[English source not present in the checked-in Bencraft/Joly files.]}
}
{
}

\Verse{123}
{
    \zA{就是这么也罢了。”}
}
{
    \eA{[English source not present in the checked-in Bencraft/Joly files.]}
}
{
}

\Verse{124}
{
    \zA{贾琏笑道：“你们太也 狠了。}
}
{
    \eA{[English source not present in the checked-in Bencraft/Joly files.]}
}
{
}

\Verse{125}
{
    \zA{你们这会子别说一千两的当头，就是现银子，要三五千，只怕也难不倒。}
}
{
    \eA{[English source not present in the checked-in Bencraft/Joly files.]}
}
{
}

\Verse{126}
{
    \zA{我 不和你们借就罢了!这会子烦你说一句话，还要个利钱，难为你们和我――”凤姐 不等说完，翻身起来说道：“我三千五千，不是赚的你的!如今里外上下，背着嚼说
        我的不少了，就短了你来说我了!可知‘没家亲引不出外鬼来’。}
}
{
    \eA{[English source not present in the checked-in Bencraft/Joly files.]}
}
{
}

\Verse{127}
{
    \zA{我们看着你家什么 石崇邓通？}
}
{
    \eA{[English source not present in the checked-in Bencraft/Joly files.]}
}
{
}

\Verse{128}
{
    \zA{把我王家的缝子扫一扫，就够你们一辈子过的了。}
}
{
    \eA{[English source not present in the checked-in Bencraft/Joly files.]}
}
{
}

\Verse{129}
{
    \zA{说出来的话也不害臊! 现有对证：把太太和我的嫁妆细看看，比一比，我们那一样是配不上你们的？”}
}
{
    \eA{[English source not present in the checked-in Bencraft/Joly files.]}
}
{
}

\Verse{130}
{
    \zA{贾 琏笑道：“说句玩话儿就急了。}
}
{
    \eA{[English source not present in the checked-in Bencraft/Joly files.]}
}
{
}

\Verse{131}
{
    \zA{这有什么的呢。}
}
{
    \eA{[English source not present in the checked-in Bencraft/Joly files.]}
}
{
}

\Verse{132}
{
    \zA{你要使一二百两银子值什么？}
}
{
    \eA{[English source not present in the checked-in Bencraft/Joly files.]}
}
{
}

\Verse{133}
{
    \zA{多的没 有，这还能够。}
}
{
    \eA{[English source not present in the checked-in Bencraft/Joly files.]}
}
{
}

\Verse{134}
{
    \zA{先拿进来，你使了再说去，如何？”}
}
{
    \eA{[English source not present in the checked-in Bencraft/Joly files.]}
}
{
}

\Verse{135}
{
    \zA{凤姐道：“我又不等着‘衔口 垫背’，忙什么呢。”}
}
{
    \eA{[English source not present in the checked-in Bencraft/Joly files.]}
}
{
}

\Verse{136}
{
    \zA{贾琏道：“何苦来？}
}
{
    \eA{[English source not present in the checked-in Bencraft/Joly files.]}
}
{
}

\Verse{137}
{
    \zA{犯不着这么肚火盛。”}
}
{
    \eA{[English source not present in the checked-in Bencraft/Joly files.]}
}
{
}

\Verse{138}
{
    \zA{凤姐听了，又笑起来， 道：“不是我着急，你说的话戳人的心。}
}
{
    \eA{[English source not present in the checked-in Bencraft/Joly files.]}
}
{
}

\Verse{139}
{
    \zA{我因为想着后日是二姐的周年，我们好了 一场，虽不能别的，到底给他上个坟，烧张纸，也是姊妹一场。}
}
{
    \eA{[English source not present in the checked-in Bencraft/Joly files.]}
}
{
}

\Verse{140}
{
    \zA{他虽没个儿女留下， 也别‘前人洒土，迷了后人的眼睛’才是。”}
}
{
    \eA{[English source not present in the checked-in Bencraft/Joly files.]}
}
{
}

\Verse{141}
{
    \zA{贾琏半晌方道：“难为你想的周全。”}
}
{
    \eA{[English source not present in the checked-in Bencraft/Joly files.]}
}
{
}

\Verse{142}
{
    \zA{凤姐一语倒把贾琏说没了话，低头打算，说：“既是后日才用，若明日得了这个， 你随便使多少就是了。”}
}
{
    \eA{[English source not present in the checked-in Bencraft/Joly files.]}
}
{
}

\Verse{143}
{
    \zA{一语未了，只见旺儿媳妇走进来。}
}
{
    \eA{[English source not present in the checked-in Bencraft/Joly files.]}
}
{
}

\Verse{144}
{
    \zA{凤姐便问：“可成了没有？”}
}
{
    \eA{[English source not present in the checked-in Bencraft/Joly files.]}
}
{
}

\Verse{145}
{
    \zA{旺儿媳妇道：“竟 不中用。}
}
{
    \eA{[English source not present in the checked-in Bencraft/Joly files.]}
}
{
}

\Verse{146}
{
    \zA{我说须得奶奶作主就成了。”}
}
{
    \eA{[English source not present in the checked-in Bencraft/Joly files.]}
}
{
}

\Verse{147}
{
    \zA{贾琏便问：“又是什么事？”}
}
{
    \eA{[English source not present in the checked-in Bencraft/Joly files.]}
}
{
}

\Verse{148}
{
    \zA{凤姐儿见问，便 说道：“不是什么大事。}
}
{
    \eA{[English source not present in the checked-in Bencraft/Joly files.]}
}
{
}

\Verse{149}
{
    \zA{旺儿有个小子，今年十七岁了还没娶媳妇儿，因要求太太 房里的彩霞，不知太太心里怎么样。}
}
{
    \eA{[English source not present in the checked-in Bencraft/Joly files.]}
}
{
}

\Verse{150}
{
    \zA{前日太太见彩霞大了，二则又多病多灾的，因 此开恩打发他出去了，给他老子随便自己择女婿去罢。}
}
{
    \eA{[English source not present in the checked-in Bencraft/Joly files.]}
}
{
}

\Verse{151}
{
    \zA{因此旺儿媳妇来求我。}
}
{
    \eA{[English source not present in the checked-in Bencraft/Joly files.]}
}
{
}

\Verse{152}
{
    \zA{我想 他两家也就算门当户对了，一说去自然成的，谁知他这会子来了说不中用。”}
}
{
    \eA{[English source not present in the checked-in Bencraft/Joly files.]}
}
{
}

\Verse{153}
{
    \zA{贾琏 道：“这是什么大事？}
}
{
    \eA{[English source not present in the checked-in Bencraft/Joly files.]}
}
{
}

\Verse{154}
{
    \zA{比彩霞好的多着呢。”}
}
{
    \eA{[English source not present in the checked-in Bencraft/Joly files.]}
}
{
}

\Verse{155}
{
    \zA{旺儿家的便笑道：“爷虽如此说，连他家 还看不起我们，别人越发看不起我们了。}
}
{
    \eA{[English source not present in the checked-in Bencraft/Joly files.]}
}
{
}

\Verse{156}
{
    \zA{好容易相看准一个媳妇儿，我只说求爷奶 奶的恩典，替作成了，奶奶又说他必是肯的，我就烦了人过去试一试，谁知白讨了 个没趣儿。}
}
{
    \eA{[English source not present in the checked-in Bencraft/Joly files.]}
}
{
}

\Verse{157}
{
    \zA{若论那孩子倒好，据我素日合意儿试他，心里没有什么说的，只是他老 子娘两个老东西太心高了些。”}
}
{
    \eA{[English source not present in the checked-in Bencraft/Joly files.]}
}
{
}

\Verse{158}
{
    \zA{一语戳动了凤姐和贾琏。}
}
{
    \eA{[English source not present in the checked-in Bencraft/Joly files.]}
}
{
}

\Verse{159}
{
    \zA{凤姐因见贾琏在此，且不做一声，只看贾琏的光景。}
}
{
    \eA{[English source not present in the checked-in Bencraft/Joly files.]}
}
{
}

\Verse{160}
{
    \zA{贾琏心中有事，那里把这点事放在心里？}
}
{
    \eA{[English source not present in the checked-in Bencraft/Joly files.]}
}
{
}

\Verse{161}
{
    \zA{待要不管，只是看着凤姐儿的陪房，且素 日出过力的，脸上实在过不去，因说：“什么大事？}
}
{
    \eA{[English source not present in the checked-in Bencraft/Joly files.]}
}
{
}

\Verse{162}
{
    \zA{只管咕咕唧唧的!你放心且去， 我明日作媒，打发两个有体面的人，一面说一面带着定礼去，就说是我的主意。}
}
{
    \eA{[English source not present in the checked-in Bencraft/Joly files.]}
}
{
}

\Verse{163}
{
    \zA{他 十分不依，叫他来见我。”}
}
{
    \eA{[English source not present in the checked-in Bencraft/Joly files.]}
}
{
}

\Verse{164}
{
    \zA{旺儿家的看着凤姐，凤姐便努嘴儿。}
}
{
    \eA{[English source not present in the checked-in Bencraft/Joly files.]}
}
{
}

\Verse{165}
{
    \zA{旺儿家的会意，忙 爬下就给贾琏磕头谢恩。}
}
{
    \eA{[English source not present in the checked-in Bencraft/Joly files.]}
}
{
}

\Verse{166}
{
    \zA{这贾琏忙道：“你只管给你们姑奶奶磕头。}
}
{
    \eA{[English source not present in the checked-in Bencraft/Joly files.]}
}
{
}

\Verse{167}
{
    \zA{我虽说了，到 底也得你们姑奶奶打发人叫他女人上来，和他好说更好些，不然太霸道了，日后你 们两亲家也难走动。”}
}
{
    \eA{[English source not present in the checked-in Bencraft/Joly files.]}
}
{
}

\Verse{168}
{
    \zA{凤姐忙道：“连你还这么开恩操心呢，我反倒袖手旁观不成？}
}
{
    \eA{[English source not present in the checked-in Bencraft/Joly files.]}
}
{
}

\Verse{169}
{
    \zA{旺儿家的你听见了：这事说了，你也忙忙的给我完了事来。}
}
{
    \eA{[English source not present in the checked-in Bencraft/Joly files.]}
}
{
}

\Verse{170}
{
    \zA{说给你男人，外头所有 的账目，一概赶今年年底都收进来，少一个钱也不依。}
}
{
    \eA{[English source not present in the checked-in Bencraft/Joly files.]}
}
{
}

\Verse{171}
{
    \zA{我的名声不好，再放一年， 都要生吃了我呢。”}
}
{
    \eA{[English source not present in the checked-in Bencraft/Joly files.]}
}
{
}

\Verse{172}
{
    \zA{旺儿媳妇笑道：“奶奶也太胆小了。}
}
{
    \eA{[English source not present in the checked-in Bencraft/Joly files.]}
}
{
}

\Verse{173}
{
    \zA{谁敢议论奶奶？}
}
{
    \eA{[English source not present in the checked-in Bencraft/Joly files.]}
}
{
}

\Verse{174}
{
    \zA{若收了时， 我也是一场痴心白使了。”}
}
{
    \eA{[English source not present in the checked-in Bencraft/Joly files.]}
}
{
}

\Verse{175}
{
    \zA{凤姐道：“我真个还等钱做什么？}
}
{
    \eA{[English source not present in the checked-in Bencraft/Joly files.]}
}
{
}

\Verse{176}
{
    \zA{不过为的是日用，出的 多，进的少。}
}
{
    \eA{[English source not present in the checked-in Bencraft/Joly files.]}
}
{
}

\Verse{177}
{
    \zA{这屋里有的没的，我和你姑爷一月的月钱，再连上四个丫头的月钱， 通共一二十两银子，还不够三五天使用的呢。}
}
{
    \eA{[English source not present in the checked-in Bencraft/Joly files.]}
}
{
}

\Verse{178}
{
    \zA{若不是我千凑万挪的，早不知过到什 么破窑里去了!如今倒落了一个放账的名儿。}
}
{
    \eA{[English source not present in the checked-in Bencraft/Joly files.]}
}
{
}

\Verse{179}
{
    \zA{既这样，我就收了回来。}
}
{
    \eA{[English source not present in the checked-in Bencraft/Joly files.]}
}
{
}

\Verse{180}
{
    \zA{我比谁不会 花钱？}
}
{
    \eA{[English source not present in the checked-in Bencraft/Joly files.]}
}
{
}

\Verse{181}
{
    \zA{咱们以后就坐着花，到多早晚就是多早晚。}
}
{
    \eA{[English source not present in the checked-in Bencraft/Joly files.]}
}
{
}

\Verse{182}
{
    \zA{这不是样儿？}
}
{
    \eA{[English source not present in the checked-in Bencraft/Joly files.]}
}
{
}

\Verse{183}
{
    \zA{前儿老太太生日，太 太急了两个月，想不出法儿来，还是我提了一句，后楼上现有些没要紧的大铜锡家 伙，四五箱子拿出去弄了三百银子，才把太太遮羞礼儿搪过去了。}
}
{
    \eA{[English source not present in the checked-in Bencraft/Joly files.]}
}
{
}

\Verse{184}
{
    \zA{我是你们知道的： 那一个金自鸣钟卖了五百六十两银子，没有半个月，大事小事没十件，白填在里头。}
}
{
    \eA{[English source not present in the checked-in Bencraft/Joly files.]}
}
{
}

\Verse{185}
{
    \zA{今儿外头也短住了，不知是谁的主意，搜寻上老太太了。}
}
{
    \eA{[English source not present in the checked-in Bencraft/Joly files.]}
}
{
}

\Verse{186}
{
    \zA{明儿再过一年，便搜寻到 头面衣裳，可就好了！”}
}
{
    \eA{[English source not present in the checked-in Bencraft/Joly files.]}
}
{
}

\Verse{187}
{
    \zA{旺儿媳妇笑道：“那一位太太奶奶的头面衣裳，折变了不够 过一辈子的？}
}
{
    \eA{[English source not present in the checked-in Bencraft/Joly files.]}
}
{
}

\Verse{188}
{
    \zA{只是不肯罢咧。”}
}
{
    \eA{[English source not present in the checked-in Bencraft/Joly files.]}
}
{
}

\Verse{189}
{
    \zA{凤姐道：“不是我说没能耐的话，要像这么着我竟不 能了。}
}
{
    \eA{[English source not present in the checked-in Bencraft/Joly files.]}
}
{
}

\Verse{190}
{
    \zA{昨儿晚上，忽然做了个梦，说来可笑：梦见一个人，虽然面善，却又不知名 姓，找我说娘娘打发他来，要一百匹锦。}
}
{
    \eA{[English source not present in the checked-in Bencraft/Joly files.]}
}
{
}

\Verse{191}
{
    \zA{我问他是那一位娘娘，他说的又不是咱们 的娘娘。}
}
{
    \eA{[English source not present in the checked-in Bencraft/Joly files.]}
}
{
}

\Verse{192}
{
    \zA{我就不肯给他，他就来夺。}
}
{
    \eA{[English source not present in the checked-in Bencraft/Joly files.]}
}
{
}

\Verse{193}
{
    \zA{正夺着，就醒了。”}
}
{
    \eA{[English source not present in the checked-in Bencraft/Joly files.]}
}
{
}

\Verse{194}
{
    \zA{旺儿家的笑道：“这是奶奶 日间操心，惦记应候宫里的事。”}
}
{
    \eA{[English source not present in the checked-in Bencraft/Joly files.]}
}
{
}

\Verse{195}
{
    \zA{一语未了，人回：“夏太监打发了一个小内家来说话。”}
}
{
    \eA{[English source not present in the checked-in Bencraft/Joly files.]}
}
{
}

\Verse{196}
{
    \zA{贾琏听了，忙皱眉道： “又是什么话？}
}
{
    \eA{[English source not present in the checked-in Bencraft/Joly files.]}
}
{
}

\Verse{197}
{
    \zA{一年他们也搬够了。”}
}
{
    \eA{[English source not present in the checked-in Bencraft/Joly files.]}
}
{
}

\Verse{198}
{
    \zA{凤姐道：“你藏起来，等我见他。}
}
{
    \eA{[English source not present in the checked-in Bencraft/Joly files.]}
}
{
}

\Verse{199}
{
    \zA{若是小事罢 了，若是大事，我自有回话。”}
}
{
    \eA{[English source not present in the checked-in Bencraft/Joly files.]}
}
{
}

\Verse{200}
{
    \zA{贾琏便躲入内套间去。}
}
{
    \eA{[English source not present in the checked-in Bencraft/Joly files.]}
}
{
}

\Verse{201}
{
    \zA{这里凤姐命人带进小太监来， 让他椅上坐了吃茶，因问何事。}
}
{
    \eA{[English source not present in the checked-in Bencraft/Joly files.]}
}
{
}

\Verse{202}
{
    \zA{那小太监便说：“夏爷爷因今儿偶见一所房子，如 今竟短二百两银子，打发我来问舅奶奶家里，有现成的银子暂借一二百，这一两日 就送来。”}
}
{
    \eA{[English source not present in the checked-in Bencraft/Joly files.]}
}
{
}

\Verse{203}
{
    \zA{凤姐儿听了，笑道：“什么是送来？}
}
{
    \eA{[English source not present in the checked-in Bencraft/Joly files.]}
}
{
}

\Verse{204}
{
    \zA{有的是银子，只管先兑了去。}
}
{
    \eA{[English source not present in the checked-in Bencraft/Joly files.]}
}
{
}

\Verse{205}
{
    \zA{改日等 我们短住，再借去也是一样。”}
}
{
    \eA{[English source not present in the checked-in Bencraft/Joly files.]}
}
{
}

\Verse{206}
{
    \zA{小太监道：“夏爷爷还说：上两回还有一千二百两银 子没送来，等今年年底下自然一齐都送过来的。”}
}
{
    \eA{[English source not present in the checked-in Bencraft/Joly files.]}
}
{
}

\Verse{207}
{
    \zA{凤姐笑道：“你夏爷爷好小气。}
}
{
    \eA{[English source not present in the checked-in Bencraft/Joly files.]}
}
{
}

\Verse{208}
{
    \zA{这 也值的放在心里？}
}
{
    \eA{[English source not present in the checked-in Bencraft/Joly files.]}
}
{
}

\Verse{209}
{
    \zA{我说一句话，不怕他多心：要都这么记清了还我们，不知要还多 少了。}
}
{
    \eA{[English source not present in the checked-in Bencraft/Joly files.]}
}
{
}

\Verse{210}
{
    \zA{只怕我们没有，要有只管拿去。”}
}
{
    \eA{[English source not present in the checked-in Bencraft/Joly files.]}
}
{
}

\Verse{211}
{
    \zA{因叫旺儿媳妇来，“出去，不管那里先支二 百银来。”}
}
{
    \eA{[English source not present in the checked-in Bencraft/Joly files.]}
}
{
}

\Verse{212}
{
    \zA{旺儿媳妇会意，因笑道：“我才因别处支不动，才来和奶奶支的。”}
}
{
    \eA{[English source not present in the checked-in Bencraft/Joly files.]}
}
{
}

\Verse{213}
{
    \zA{凤姐 道：“你们只会里头来要钱，叫你们外头弄去，就不能了。”}
}
{
    \eA{[English source not present in the checked-in Bencraft/Joly files.]}
}
{
}

\Verse{214}
{
    \zA{说着，叫平儿：“把我 那两个金项圈拿出去，暂且押四百两银子。”}
}
{
    \eA{[English source not present in the checked-in Bencraft/Joly files.]}
}
{
}

\Verse{215}
{
    \zA{平儿答应去了，果然拿了一个锦盒子 来，里面两个锦袱包着。}
}
{
    \eA{[English source not present in the checked-in Bencraft/Joly files.]}
}
{
}

\Verse{216}
{
    \zA{打开时，一个金累丝攒珠的，那珍珠都有莲子大小；一个 点翠嵌宝石的：两个都与宫中之物不离上下。}
}
{
    \eA{[English source not present in the checked-in Bencraft/Joly files.]}
}
{
}

\Verse{217}
{
    \zA{一时拿去，果然拿了四百两银子来。}
}
{
    \eA{[English source not present in the checked-in Bencraft/Joly files.]}
}
{
}

\Verse{218}
{
    \zA{凤姐命给小太监打叠一半，那一半与了旺儿媳妇，命他拿去办八月中秋的节。}
}
{
    \eA{[English source not present in the checked-in Bencraft/Joly files.]}
}
{
}

\Verse{219}
{
    \zA{那小 太监便告辞了，凤姐命人替他拿着银子，送出大门去了。}
}
{
    \eA{[English source not present in the checked-in Bencraft/Joly files.]}
}
{
}

\Verse{220}
{
    \zA{这里贾琏出来笑道：“这 一起外祟，何日是了！”}
}
{
    \eA{[English source not present in the checked-in Bencraft/Joly files.]}
}
{
}

\Verse{221}
{
    \zA{凤姐笑道：“刚说着，就来了一股子。”}
}
{
    \eA{[English source not present in the checked-in Bencraft/Joly files.]}
}
{
}

\Verse{222}
{
    \zA{贾琏道：“昨儿周太 监来，张口一千两，我略应慢了些，他就不自在。}
}
{
    \eA{[English source not present in the checked-in Bencraft/Joly files.]}
}
{
}

\Verse{223}
{
    \zA{将来得罪人的地方儿多着呢。}
}
{
    \eA{[English source not present in the checked-in Bencraft/Joly files.]}
}
{
}

\Verse{224}
{
    \zA{这 会子再发个三五万的财就好了！”}
}
{
    \eA{[English source not present in the checked-in Bencraft/Joly files.]}
}
{
}

\Verse{225}
{
    \zA{一面说，一面平儿伏侍凤姐另洗了脸、更衣，往 贾母处伺候晚饭。}
}
{
    \eA{[English source not present in the checked-in Bencraft/Joly files.]}
}
{
}

\Verse{226}
{
    \zA{这里贾琏出来，刚至外书房，忽见林之孝走来。}
}
{
    \eA{[English source not present in the checked-in Bencraft/Joly files.]}
}
{
}

\Verse{227}
{
    \zA{贾琏因问何事。}
}
{
    \eA{[English source not present in the checked-in Bencraft/Joly files.]}
}
{
}

\Verse{228}
{
    \zA{林之孝说道： “才听见雨村降了，却不知何事。}
}
{
    \eA{[English source not present in the checked-in Bencraft/Joly files.]}
}
{
}

\Verse{229}
{
    \zA{只怕未必真。”}
}
{
    \eA{[English source not present in the checked-in Bencraft/Joly files.]}
}
{
}

\Verse{230}
{
    \zA{贾琏道：“真不真，他那官儿未必 保的长。}
}
{
    \eA{[English source not present in the checked-in Bencraft/Joly files.]}
}
{
}

\Verse{231}
{
    \zA{只怕将来有事，咱们宁可疏远着他好。”}
}
{
    \eA{[English source not present in the checked-in Bencraft/Joly files.]}
}
{
}

\Verse{232}
{
    \zA{林之孝道：“何从不是？}
}
{
    \eA{[English source not present in the checked-in Bencraft/Joly files.]}
}
{
}

\Verse{233}
{
    \zA{只是一时 难以疏远。}
}
{
    \eA{[English source not present in the checked-in Bencraft/Joly files.]}
}
{
}

\Verse{234}
{
    \zA{如今东府大爷和他更好，老爷又喜欢他，时常来往，那个不知？”}
}
{
    \eA{[English source not present in the checked-in Bencraft/Joly files.]}
}
{
}

\Verse{235}
{
    \zA{贾琏 道：“横竖不和他谋事，也不相干。}
}
{
    \eA{[English source not present in the checked-in Bencraft/Joly files.]}
}
{
}

\Verse{236}
{
    \zA{你去再打听真了是为什么。”}
}
{
    \eA{[English source not present in the checked-in Bencraft/Joly files.]}
}
{
}

\Verse{237}
{
    \zA{林之孝答应了，却 不动身，坐在椅子上再说闲话。}
}
{
    \eA{[English source not present in the checked-in Bencraft/Joly files.]}
}
{
}

\Verse{238}
{
    \zA{因又说起家道艰难，便趁势说：“人口太众了。}
}
{
    \eA{[English source not present in the checked-in Bencraft/Joly files.]}
}
{
}

\Verse{239}
{
    \zA{不 如拣个空日回明老太太老爷，把这些出过力的老家人，用不着的，开恩放几家出去： 一则他们各有营运，二则家里一年也省口粮月钱。}
}
{
    \eA{[English source not present in the checked-in Bencraft/Joly files.]}
}
{
}

\Verse{240}
{
    \zA{再者，里头的姑娘也太多。}
}
{
    \eA{[English source not present in the checked-in Bencraft/Joly files.]}
}
{
}

\Verse{241}
{
    \zA{俗语 说，‘一时比不得一时’如今说不得先时的例了，少不的大家委屈些，该使八个的 使六个，使四个的使两个。}
}
{
    \eA{[English source not present in the checked-in Bencraft/Joly files.]}
}
{
}

\Verse{242}
{
    \zA{若各房算起来，一年也可以省许多月米月钱。}
}
{
    \eA{[English source not present in the checked-in Bencraft/Joly files.]}
}
{
}

\Verse{243}
{
    \zA{况且里头 的女孩子们，一半都大了，也该配人的配人，成了房，岂不又滋生出些人来？”}
}
{
    \eA{[English source not present in the checked-in Bencraft/Joly files.]}
}
{
}

\Verse{244}
{
    \zA{贾 琏道：“我也这么想，只是老爷才回家来，多少大事未回，那里议到这个上头？}
}
{
    \eA{[English source not present in the checked-in Bencraft/Joly files.]}
}
{
}

\Verse{245}
{
    \zA{前儿 官媒拿了个庚帖来求亲，太太还说老爷才来家，每日欢天喜地的说骨肉完聚，忽然 提起这事，恐老爷又伤心，所以且不叫提起。”}
}
{
    \eA{[English source not present in the checked-in Bencraft/Joly files.]}
}
{
}

\Verse{246}
{
    \zA{林之孝道：“这也是正理，太太想的 周到。”}
}
{
    \eA{[English source not present in the checked-in Bencraft/Joly files.]}
}
{
}

\Verse{247}
{
    \zA{贾琏道：“正是，提起这话，我想起一件事来：我们旺儿的小子，要说太太 屋里的彩霞，他昨儿求我。}
}
{
    \eA{[English source not present in the checked-in Bencraft/Joly files.]}
}
{
}

\Verse{248}
{
    \zA{我想什么大事，不管谁去说一声去，就说我的话。”}
}
{
    \eA{[English source not present in the checked-in Bencraft/Joly files.]}
}
{
}

\Verse{249}
{
    \zA{林 之孝答应了，半晌笑道：“依我说，二爷竟别管这件事。}
}
{
    \eA{[English source not present in the checked-in Bencraft/Joly files.]}
}
{
}

\Verse{250}
{
    \zA{旺儿的那小子虽然年轻， 在外吃酒赌钱，无所不至。}
}
{
    \eA{[English source not present in the checked-in Bencraft/Joly files.]}
}
{
}

\Verse{251}
{
    \zA{虽说都是奴才，到底是一辈子的事。}
}
{
    \eA{[English source not present in the checked-in Bencraft/Joly files.]}
}
{
}

\Verse{252}
{
    \zA{彩霞这孩子这几年 我虽没看见，听见说越发出跳的好了，何苦来白遭塌一个人呢？”}
}
{
    \eA{[English source not present in the checked-in Bencraft/Joly files.]}
}
{
}

\Verse{253}
{
    \zA{贾琏道：“哦!他 小子竟会喝酒不成人吗？}
}
{
    \eA{[English source not present in the checked-in Bencraft/Joly files.]}
}
{
}

\Verse{254}
{
    \zA{这么着，那里还给他老婆？}
}
{
    \eA{[English source not present in the checked-in Bencraft/Joly files.]}
}
{
}

\Verse{255}
{
    \zA{且给他一顿棍，锁起来，再问他 老子娘。”}
}
{
    \eA{[English source not present in the checked-in Bencraft/Joly files.]}
}
{
}

\Verse{256}
{
    \zA{林之孝笑道：“何必在这一时？}
}
{
    \eA{[English source not present in the checked-in Bencraft/Joly files.]}
}
{
}

\Verse{257}
{
    \zA{等他再生事，我们自然回爷处治，如今且 也不用究办。”}
}
{
    \eA{[English source not present in the checked-in Bencraft/Joly files.]}
}
{
}

\Verse{258}
{
    \zA{贾琏不语。}
}
{
    \eA{[English source not present in the checked-in Bencraft/Joly files.]}
}
{
}

\Verse{259}
{
    \zA{一时林之孝出去。}
}
{
    \eA{[English source not present in the checked-in Bencraft/Joly files.]}
}
{
}

\Verse{260}
{
    \zA{晚间凤姐已命人唤了彩霞之母来说媒。}
}
{
    \eA{[English source not present in the checked-in Bencraft/Joly files.]}
}
{
}

\Verse{261}
{
    \zA{那彩霞之母满心纵不愿意，见凤姐自和 他说，何等体面，便心不由己的满口应了出去。}
}
{
    \eA{[English source not present in the checked-in Bencraft/Joly files.]}
}
{
}

\Verse{262}
{
    \zA{凤姐又问贾琏：“可说了没有？”}
}
{
    \eA{[English source not present in the checked-in Bencraft/Joly files.]}
}
{
}

\Verse{263}
{
    \zA{贾琏因说：“我原要说来着，听见他这小子大不成人，所以还没说。}
}
{
    \eA{[English source not present in the checked-in Bencraft/Joly files.]}
}
{
}

\Verse{264}
{
    \zA{若果然不成人， 且管教他两日，再给他老婆不迟。”}
}
{
    \eA{[English source not present in the checked-in Bencraft/Joly files.]}
}
{
}

\Verse{265}
{
    \zA{凤姐笑道：“我们王家的人，连我还不中你们的 意，何况奴才呢。}
}
{
    \eA{[English source not present in the checked-in Bencraft/Joly files.]}
}
{
}

\Verse{266}
{
    \zA{我已经和他娘说了，他娘倒欢天喜地，难道又叫进他来不要了不 成？”}
}
{
    \eA{[English source not present in the checked-in Bencraft/Joly files.]}
}
{
}

\Verse{267}
{
    \zA{贾琏道：“你既说了，又何必退呢？}
}
{
    \eA{[English source not present in the checked-in Bencraft/Joly files.]}
}
{
}

\Verse{268}
{
    \zA{明日说给他老子，好生管他就是了。”}
}
{
    \eA{[English source not present in the checked-in Bencraft/Joly files.]}
}
{
}

\Verse{269}
{
    \zA{这 里说话不提。}
}
{
    \eA{[English source not present in the checked-in Bencraft/Joly files.]}
}
{
}

\Verse{270}
{
    \zA{且说彩霞因前日出去等父母择人，心中虽与贾环有旧，尚未作准。}
}
{
    \eA{[English source not present in the checked-in Bencraft/Joly files.]}
}
{
}

\Verse{271}
{
    \zA{今日又见旺 儿每每来求亲，早闻得旺儿之子酗酒赌博，而且容颜丑陋，不能如意。}
}
{
    \eA{[English source not present in the checked-in Bencraft/Joly files.]}
}
{
}

\Verse{272}
{
    \zA{自此，心中 越发懊恼，惟恐旺儿仗势作成，终身不遂，未免心中急躁。}
}
{
    \eA{[English source not present in the checked-in Bencraft/Joly files.]}
}
{
}

\Verse{273}
{
    \zA{至晚间，悄命他妹子小 霞进二门来找赵姨娘，问个端底。}
}
{
    \eA{[English source not present in the checked-in Bencraft/Joly files.]}
}
{
}

\Verse{274}
{
    \zA{赵姨娘素日深与彩霞好，巴不得给了贾环，方有 个膀臂，不承望王夫人又放出去了。}
}
{
    \eA{[English source not present in the checked-in Bencraft/Joly files.]}
}
{
}

\Verse{275}
{
    \zA{每每调唆贾环去讨，一则贾环羞口难开，二则 贾环也不在意，不过是个丫头，他去了将来自然还有好的，遂迁延住不肯说去，意 思便丢开了手。}
}
{
    \eA{[English source not present in the checked-in Bencraft/Joly files.]}
}
{
}

\Verse{276}
{
    \zA{无奈赵姨娘又不舍，又见他妹子来问，是晚得空，便先求了贾政。}
}
{
    \eA{[English source not present in the checked-in Bencraft/Joly files.]}
}
{
}

\Verse{277}
{
    \zA{贾政说道：“且忙什么。}
}
{
    \eA{[English source not present in the checked-in Bencraft/Joly files.]}
}
{
}

\Verse{278}
{
    \zA{等他们再念一二年书，再放人不迟。}
}
{
    \eA{[English source not present in the checked-in Bencraft/Joly files.]}
}
{
}

\Verse{279}
{
    \zA{我已经看中了两个丫 头，一个给宝玉，一个给环儿。}
}
{
    \eA{[English source not present in the checked-in Bencraft/Joly files.]}
}
{
}

\Verse{280}
{
    \zA{只是年纪还小，又怕他们误了念书，再等一二年再 提。”}
}
{
    \eA{[English source not present in the checked-in Bencraft/Joly files.]}
}
{
}

\Verse{281}
{
    \zA{赵姨娘还要说话，只听外面一声响，不知何物，大家吃了一惊。}
}
{
    \eA{[English source not present in the checked-in Bencraft/Joly files.]}
}
{
}

\Verse{282}
{
    \zA{未知如何，下回分解。}
}
{
    \eA{[English source not present in the checked-in Bencraft/Joly files.]}
}
{
}

\Verse{283}
{
    \zA{(本章完)}
}
{
    \eA{[English source not present in the checked-in Bencraft/Joly files.]}
}
{
}
